% Einheit 1 von 4 – Gleichungen verstehen (bank/lineare-gleichungen/e1.jsonl, zusammenbau v0.1)
%% TODO Zweigzeile Teil 1: was hier gelernt wird (Satz in Schülersprache aus dem Einheitstitel) – Einheit 1
\einheitenkopf[e1]{Einheit 1 von 4 · Gleichungen verstehen}
\zweigzeile{ab Kl. 5, je nach Buch bis Kl. 6 (am Gymnasium ab Kl. 6) · P10}
\setcounter{aufgabe}{7}

%% TODO Titel: Ich-kann-Satz fehlt in der Bank (Platzhalter: Kettenname) – Nr. 8
%% TODO Anweisung: ein Satz über der ersten Teilaufgabe fehlt in der Bank – Nr. 8
\begin{aufgabe}{Umkehroperation benennen}
\begin{teile}
\teil Was macht $+\,6$ rückgängig? \leerfeld
\end{teile}
\end{aufgabe}

%% TODO \verfahren{…}: Verfahrensüberschrift in Sachsprache fehlt in der Bank (2.3 b)
%% TODO Titel: Ich-kann-Satz fehlt in der Bank (Platzhalter: Kettenname) – Nr. 9
%% TODO Anweisung: ein Satz über der ersten Teilaufgabe fehlt in der Bank – Nr. 9
\begin{aufgabe}{Lösung prüfen}
%% TODO \rechenplatz{n}: Schrittzahl des Grundfalls fehlt in der Bank (nur bei fester Schreibform, 2.3 b)
\begin{teile}
\teil $x + 9 = 15$ – ist $6$ Lösung?
\teil $2x + 7 = 13$ – ist $3$ Lösung?
\teil $3x + 4 = 25$ – ist $8$ Lösung?
\teil $5x - 2 = 3x + 6$ – ist $4$ Lösung?
\teil Welche Gleichung hat die Lösung $-3$? \hfill (P10 2018 OS) \\ \kreuz{$2x - 6 = 0$} \\ \kreuz{$-4x + 12 = 0$} \\ \kreuz{$3x + 5 = -4$} \\ \kreuz{$6x + 20 = 1$}
\end{teile}
\end{aufgabe}

%% TODO \verfahren{…}: Verfahrensüberschrift in Sachsprache fehlt in der Bank (2.3 b)
%% TODO Titel: Ich-kann-Satz fehlt in der Bank (Platzhalter: Kettenname) – Nr. 10
%% TODO Anweisung: ein Satz über der ersten Teilaufgabe fehlt in der Bank – Nr. 10
\begin{aufgabe}{Durch Probieren lösen}
%% TODO \rechenplatz{n}: Schrittzahl des Grundfalls fehlt in der Bank (nur bei fester Schreibform, 2.3 b)
\begin{teile}
\teil Für welches $x$ ist $2x + 3 = 11$?

\wertetabelle{x}{$2x + 3$}{1,2,3,4}
x = \leerfeld
\teil Probiere eigene Zahlen: $3x + 8 = 38$

\wertetabelleleer{x}{$3x + 8$}{4}
x = \leerfeld
\teil $x + 13 = 21$ – welche Zahl plus 13 ergibt 21? x = \leerfeld
\teil $7x - 6 = 22$ – welche der Zahlen 3, 4, 5 ist die Lösung? Warum?
\teil $x + 47 = 83$ – probiere in einer Tabelle, nenne die Lösung und mache die Probe.

\wertetabelleleer{x}{$x + 47$}{4}
x = \leerfeld
\end{teile}
\end{aufgabe}

%% TODO Titel: Ich-kann-Satz fehlt in der Bank (Platzhalter: Kettenname) – Nr. 11
%% TODO Anweisung: ein Satz über der ersten Teilaufgabe fehlt in der Bank – Nr. 11
\begin{aufgabe}{Lösung prüfen – Fehler finden, Begründen, Darstellungswechsel}
\begin{teile}
\teil Ben prüft, ob 4 Lösung von $3x + 2 = x + 10$ ist. \rechnung{3 \cdot 4 + 2 &= 14 \\ 14 &\ne 10 && \text{(fA)}} Finde den Fehler.
\teil Mira löst $4x + 6 = 30$ und erhält 9. Eine Probe macht sie nicht. Finde den Fehler.
\teil Kann $x + 25 = 10$ eine positive Lösung haben? Begründe ohne Rechnung.
\teil „Welche Zahl minus 6 ergibt 15?" – als Gleichung
\teil $x + 14 = 30$ – als Frage „Welche Zahl …?" und Antwort
\end{teile}
\end{aufgabe}
